\documentclass[10pt]{article}

\usepackage[T1]{fontenc}
\usepackage[utf8]{inputenc}
\usepackage{lmodern}
\usepackage{microtype}
\usepackage[margin=1in]{geometry}
\usepackage{amsmath,amssymb}
\usepackage{booktabs}
\usepackage{multirow}
\usepackage{array}
\usepackage{graphicx}
\usepackage{xcolor}
\usepackage{enumitem}
\usepackage{hyperref}
\usepackage[nameinlink,capitalize]{cleveref}

\hypersetup{
  colorlinks=true,
  linkcolor=blue!55!black,
  citecolor=green!45!black,
  urlcolor=blue!65!black
}

\newcommand{\method}{MI-VLN}
\newcommand{\tbd}[1]{\textcolor{red!75!black}{[TBD: #1]}}

\title{\textbf{Ask or Intervene: A Benchmark for\\
Mixed-Initiative Vision-and-Language Navigation}}
\author{Anonymous Authors}
\date{}

\begin{document}
\maketitle

\begin{abstract}
Vision-and-language navigation (VLN) typically assumes that a robot receives
a fixed instruction before execution and then navigates autonomously.
Real service robots instead operate under evolving human intent: a user may
correct a mistaken turn, add a constraint, replace the destination, or answer
a robot's request for assistance. Importantly, the user need not continuously
monitor the robot, but must be able to intervene whenever necessary.
We introduce \method, a benchmark for \emph{mixed-initiative} VLN in indoor
service environments. The benchmark models two complementary interaction
directions: robot-initiated help requests and human-initiated interventions.
It evaluates three tracks---human-initiated, robot-initiated, and fully
mixed-initiative navigation---under asynchronous user availability.
To isolate the value of interaction, each episode contains a shared pre-event
trajectory and counterfactual post-event branches with matched initial
conditions. We define intervention windows, dynamic task specifications,
interaction-aware metrics, and a reproducible annotation protocol based on
Habitat 3.0 and HSSD scenes. We further propose a baseline agent that combines
visual-spatial memory, dynamic task memory, an initiative manager, and a
language-conditioned navigator. Experiments are designed to measure navigation
success, recovery efficiency, communication cost, intervention timing, and
robustness to delayed human responses. \tbd{Insert the main quantitative
finding after completing experiments.}
\end{abstract}

\section{Introduction}
\label{sec:intro}

An instruction such as ``go to the kitchen and stop beside the table'' is
rarely immutable in a real home. While the robot is moving, a user may notice
that the intended table was ambiguous, ask the robot to avoid a wet floor, or
change the destination because the household task has changed. Conversely, the
robot may recognize that its observation is ambiguous and ask whether it
should take the left or right corridor. These interactions form a continuous
decision process in which control of the dialogue can shift between the human
and the robot.

Most VLN benchmarks evaluate a different setting: an agent receives a fixed
natural-language instruction and is judged primarily by whether it reaches a
static goal \cite{anderson2018r2r,ku2020rxr,krantz2020beyond}. Interactive VLN
work has studied dialogue-based assistance or robot help-seeking
\cite{thomason2020cvdn,nguyen2019hanna,zhu2020cmn}, while online robot
correction has studied how language modifies plans or policies
\cite{sharma2022correcting,broad2017language}. However, a practical service
robot must support both directions of initiative in the same episode. It must
decide when to act, when to ask, how to incorporate unsolicited corrections,
and how to proceed when a human is temporarily unavailable.

This paper introduces \method, a benchmark that treats interaction as part of
the navigation problem rather than as an external reset mechanism. A human
may intervene at any point, but is not assumed to watch every step. Likewise,
the robot may request assistance, but unnecessary questions incur a cost.
Messages update a structured task state without erasing the agent's navigation
history. The benchmark therefore exposes a central trade-off: autonomous
progress versus timely, useful communication.

Our contributions are:
\begin{itemize}[leftmargin=1.4em]
  \item We formulate mixed-initiative VLN with bidirectional initiative,
  asynchronous human availability, and a dynamic language-defined task state.
  \item We design a household-service benchmark with three evaluation tracks,
  controlled intervention windows, counterfactual branches, and event-level
  annotations.
  \item We introduce metrics that jointly evaluate task success, recovery,
  communication cost, timing quality, and human effort.
  \item We provide a modular baseline and an experimental protocol for
  studying when an agent should ask, how it should adapt, and when human
  intervention is most valuable.
\end{itemize}

\begin{figure}[t]
\centering
\fbox{\begin{minipage}{0.94\linewidth}
\small
\textbf{Episode timeline.}
The robot receives an initial instruction and navigates autonomously.
At an event window, either (a) the human intervenes with a correction or a
new constraint, or (b) the robot asks a targeted question and may receive a
delayed answer. The interaction updates the current task specification; the
physical state and navigation history are preserved. Evaluation continues
until success, failure, timeout, or explicit termination.
\end{minipage}}
\caption{Mixed-initiative navigation as a continuous, interruptible process.
The benchmark does not require continuous human observation.}
\label{fig:overview}
\end{figure}

\section{Related Work}
\label{sec:related}

\paragraph{Vision-and-language navigation.}
R2R established instruction-guided navigation in photo-realistic indoor
environments \cite{anderson2018r2r}. Subsequent benchmarks expanded language,
path diversity, and continuous control
\cite{ku2020rxr,krantz2020beyond,krantz2023iterative}. These settings usually
treat the instruction and goal as fixed throughout an episode. \method instead
evaluates how an agent responds when intent changes during execution.

\paragraph{Dialogue and help-seeking in navigation.}
CVDN studies navigation through question--answer dialogue
\cite{thomason2020cvdn}. HANNA allows an agent to request help from an
assistant \cite{nguyen2019hanna}, and later work studies help-request
generation and collaborative navigation
\cite{zhu2020cmn,zhu2021scoa,taioli2025collaborative}. These works motivate the
robot-initiated component of \method. Our benchmark adds unsolicited
human-initiated intervention, temporary human unavailability, and direct
comparison between initiative policies.

\paragraph{Online correction and shared autonomy.}
Language can update a robot policy or repair an erroneous plan during
execution \cite{broad2017language,sharma2022correcting,cui2023lilac}.
Shared-autonomy research further emphasizes adjustable control and human effort
\cite{beer2014levels}. We translate these ideas to embodied VLN, where the
agent must ground an evolving task specification in egocentric observations
and long-horizon spatial memory.

\paragraph{Embodied simulation.}
Habitat provides a scalable platform for embodied AI
\cite{savva2019habitat}; Habitat 3.0 extends it toward collaborative
human--robot tasks \cite{puig2024habitat3}. HSSD supplies high-quality,
semantically annotated residential scenes \cite{khanna2024hssd}. We use these
assets to construct reproducible household-service episodes.

\section{Task Formulation}
\label{sec:task}

\subsection{Interaction Model}

At time $t$, the agent receives an egocentric observation
\begin{equation}
  o_t = (I_t^{\mathrm{rgb}}, I_t^{\mathrm{depth}}, p_t),
\end{equation}
where $p_t$ denotes pose information when available. The interaction history is
$H_t=\{u_0,m_1,\ldots,m_t\}$, where $u_0$ is the initial instruction and each
$m_i$ is an optional human or robot message.

The agent chooses from navigation and communication actions:
\begin{equation}
  a_t \in \mathcal{A}_{\mathrm{nav}} \cup \mathcal{A}_{\mathrm{ask}},
\end{equation}
where $\mathcal{A}_{\mathrm{nav}}$ contains movement and stop actions, and
$\mathcal{A}_{\mathrm{ask}}$ contains typed help requests. A human message may
be a correction, route constraint, subgoal insertion, goal replacement,
clarification, confirmation, or abort command.

Human availability is represented by $h_t\in\{0,1\}$. If the agent asks when
$h_t=0$, the question may remain unanswered, be answered after a sampled
delay, or expire according to the evaluation condition. This prevents the
benchmark from equating ``intervenable'' with ``continuously supervised.''

\subsection{Dynamic Task State}

Instead of repeatedly concatenating all dialogue, we maintain a structured
task state
\begin{equation}
  B_t=(g_t,\mathcal{S}_t,\mathcal{C}_t,\mathcal{R}_t,q_t),
\end{equation}
where $g_t$ is the current goal, $\mathcal{S}_t$ is an ordered subgoal list,
$\mathcal{C}_t$ contains route and safety constraints, $\mathcal{R}_t$
contains resolved references, and $q_t$ is the current interaction status.
A message induces an update
\begin{equation}
  B_t=f_{\mathrm{update}}(B_{t-1},m_t).
\end{equation}
Crucially, this update does not reset the physical state or erase visual and
spatial memory.

\subsection{Evaluation Tracks}

\begin{enumerate}[leftmargin=1.6em]
  \item \textbf{Human-Initiated (HI):} the robot cannot ask questions, but
  must process interventions arriving during navigation.
  \item \textbf{Robot-Initiated (RI):} the task is fixed unless the robot asks
  for assistance; asking decisions and question utility are evaluated.
  \item \textbf{Mixed-Initiative (MI):} both parties may initiate interaction,
  possibly under delayed or missing human responses.
\end{enumerate}

The agent seeks to maximize task utility
\begin{equation}
  J=R_{\mathrm{task}}
  -\lambda_{\mathrm{path}}C_{\mathrm{path}}
  -\lambda_{\mathrm{comm}}C_{\mathrm{comm}}
  -\lambda_{\mathrm{risk}}C_{\mathrm{risk}},
\end{equation}
rather than maximizing success at the cost of unlimited human assistance.

\section{The \method\ Benchmark}
\label{sec:benchmark}

\subsection{Environment and Embodiment}

The benchmark uses Habitat 3.0 with HSSD residential scenes. The primary
embodiment is a single mobile service robot with RGB-D sensing and discrete
short-horizon actions: move forward, turn left, turn right, look up, look down,
and stop. This constrained action space makes intervention effects easier to
replay and compare. A continuous-control extension is left for future work.

\subsection{Base Household Tasks}

Episodes cover navigation patterns common in domestic service:
\begin{itemize}[leftmargin=1.4em]
  \item \textbf{Room and object approach:} reach a kitchen counter, bedside
  table, sofa, charging station, or storage shelf.
  \item \textbf{Person-oriented service:} navigate to a resident's location,
  wait nearby, or return to the user after visiting another room.
  \item \textbf{Multi-stage service:} visit a source location and then a
  delivery location, represented initially as ordered navigation subgoals.
  \item \textbf{Constraint-aware routing:} avoid a room, doorway, fragile
  area, occupied space, or temporarily blocked corridor.
  \item \textbf{Dynamic intent:} replace the destination, reorder subgoals,
  cancel an obsolete task, or redirect the robot to an urgent location.
\end{itemize}
The first release focuses on navigation and does not require physical object
manipulation, although task descriptions may supply realistic service context.

\subsection{Intervention Taxonomy}

Each intervention is labeled by its semantic operation:
\begin{table}[t]
\centering
\small
\begin{tabular}{p{0.25\linewidth}p{0.62\linewidth}}
\toprule
\textbf{Type} & \textbf{Example and state change}\\
\midrule
Directional correction &
``Not that hallway; turn back and take the door on your left.''\\
Reference repair &
``By the table, I meant the small table beside the sofa.''\\
Constraint addition &
``Do not enter the nursery; the baby is sleeping.''\\
Subgoal insertion &
``Before the kitchen, stop at the charging dock.''\\
Goal replacement &
``Forget the bedroom; come back to the living room.''\\
Priority or abort &
``Stop there'' or ``Return immediately.''\\
\bottomrule
\end{tabular}
\caption{Human-initiated intervention types in the initial benchmark.}
\label{tab:interventions}
\end{table}

Robot questions are similarly typed as direction, landmark, reference,
constraint, goal, or confirmation queries. Templates may be used for controlled
training data, while test messages include human paraphrases.

\subsection{Selecting Intervention Times}
\label{sec:timing}

Interventions should be consequential but not arbitrary. We define candidate
windows from observable episode events:
\begin{enumerate}[leftmargin=1.6em]
  \item \textbf{Pre-error ambiguity:} before a junction with two plausible
  routes or referents.
  \item \textbf{Recoverable deviation:} after the agent commits to an
  incorrect route but before recovery becomes trivial or impossible.
  \item \textbf{Constraint relevance:} shortly before entering a newly
  prohibited region.
  \item \textbf{Intent change:} after meaningful progress, when a replacement
  goal produces a nontrivial replanning problem.
  \item \textbf{Safety-critical boundary:} before collision-prone or
  human-occupied areas, subject to conservative simulator rules.
\end{enumerate}

For each window, we sample early, middle, and late intervention points. Let
$\Delta C(t)$ denote the additional optimal path cost caused by waiting until
$t$ to intervene. Episodes are retained only when intervention changes the
optimal continuation by at least threshold $\epsilon$ and still permits task
completion. This yields meaningful timing comparisons without selecting times
from a particular learned agent's failures alone.

\subsection{Counterfactual Episode Branches}

Each benchmark item stores a shared prefix state $s_\tau$ and several branches:
\begin{equation}
  \mathcal{E}=
  \left(\xi_{0:\tau},e_\tau,
  \xi_{\tau:T}^{\mathrm{no\text{-}int}},
  \xi_{\tau:T}^{\mathrm{int}},
  \xi_{\tau:T}^{\mathrm{oracle}}\right),
\end{equation}
where $e_\tau$ is an intervention opportunity. All branches begin from the
same simulator state. This design supports paired evaluation of no
interaction, actual interaction, and oracle-timed interaction, reducing
variance caused by unrelated trajectory differences.

\subsection{Annotation and Quality Control}

Annotation is performed through a Windows client connected to a Linux Habitat
backend. Annotators receive a short observation clip, a top-down map available
only to the annotation interface, the current instruction, and the intended
state update. They produce a natural intervention or response and label its
semantic operation. A second annotator verifies groundedness, executability,
and consistency with the intended update. Paraphrase diversity is measured
with lexical and semantic similarity statistics; contradictory or
under-specified messages are adjudicated.

\subsection{Data Splits}

We partition by scene into train, validation-seen, validation-unseen, and
test-unseen sets. Test splits additionally hold out selected combinations of
intervention type, timing, and room transition to evaluate compositional
generalization. No scene geometry or human-authored paraphrase is shared
between train and unseen test sets.

\begin{table}[t]
\centering
\small
\begin{tabular}{lrrrr}
\toprule
\textbf{Split} & \textbf{Scenes} & \textbf{Base eps.} &
\textbf{Branches} & \textbf{Utterances}\\
\midrule
Train & \tbd{N} & \tbd{N} & \tbd{N} & \tbd{N}\\
Val-Seen & \tbd{N} & \tbd{N} & \tbd{N} & \tbd{N}\\
Val-Unseen & \tbd{N} & \tbd{N} & \tbd{N} & \tbd{N}\\
Test-Unseen & \tbd{N} & \tbd{N} & \tbd{N} & \tbd{N}\\
\bottomrule
\end{tabular}
\caption{Planned benchmark statistics. Replace placeholders after generation
and quality control.}
\label{tab:data}
\end{table}

\subsection{Metrics}

We report standard navigation success rate (SR), success weighted by path
length (SPL), and navigation error. Interaction is measured by:
\begin{align}
  \mathrm{RecoveryRate}&=
  \frac{N_{\mathrm{failed\ without\ interaction,\ succeeded\ with\ interaction}}}
       {N_{\mathrm{recoverable\ episodes}}},\\
  \mathrm{AskPrecision}&=
  \frac{N_{\mathrm{useful\ questions}}}{N_{\mathrm{all\ questions}}},\\
  \mathrm{AskRecall}&=
  \frac{N_{\mathrm{help\ needed\ and\ requested}}}
       {N_{\mathrm{help\ needed}}},\\
  \mathrm{UpdateConsistency}&=
  \frac{N_{\mathrm{actions\ consistent\ with\ latest\ task}}}
       {N_{\mathrm{post\text{-}message\ actions}}}.
\end{align}
We also report question count, human intervention count, response delay,
distance traveled after obsolete goals, constraint violations, and success
gain per human message. A combined score is provided only as a secondary
summary because its cost weights reflect deployment preferences.

\section{A Mixed-Initiative Baseline}
\label{sec:method}

Our baseline contains four modules.

\paragraph{Visual-spatial memory.}
An RGB-D encoder produces observation features that are accumulated in a
topological or metric memory $M_t^{\mathrm{map}}$. This memory records visited
locations, visual landmarks, and traversability evidence.

\paragraph{Dynamic task memory.}
A language encoder maps each new message to a semantic update over $B_t$.
The update module predicts the affected fields and whether prior goals or
constraints are superseded. The resulting memory preserves both the current
task and an auditable change history.

\paragraph{Initiative manager.}
The manager estimates uncertainty $U_t$, expected failure risk $F_t$, and the
expected value of a question $V_t^{\mathrm{ask}}$. It asks only if
\begin{equation}
  V_t^{\mathrm{ask}}>C_{\mathrm{ask}}+C_{\mathrm{delay}},
\end{equation}
otherwise it navigates or waits. During human unavailability, it may continue
with a conservative action, pause at a safe location, or postpone the
question.

\paragraph{Navigator.}
A policy conditions on the current observation, spatial memory, task state,
and interaction status:
\begin{equation}
  \pi(a_t\mid o_t,M_t^{\mathrm{map}},B_t,H_t).
\end{equation}
Training combines imitation learning on expert branches with reinforcement
learning or offline preference optimization for communication cost.
The total loss is
\begin{equation}
  \mathcal{L}=\mathcal{L}_{\mathrm{nav}}
  +\alpha\mathcal{L}_{\mathrm{update}}
  +\beta\mathcal{L}_{\mathrm{ask}}
  +\gamma\mathcal{L}_{\mathrm{consistency}}.
\end{equation}

\section{Experiments}
\label{sec:experiments}

\subsection{Research Questions}

We organize experiments around five questions:
\begin{enumerate}[leftmargin=1.6em]
  \item Does mid-execution language improve task success and recovery?
  \item Can an agent identify when asking is more valuable than acting?
  \item How do human- and robot-initiated interactions complement each other?
  \item How sensitive are agents to intervention timing and delayed responses?
  \item Does structured task memory improve adaptation over dialogue
  concatenation?
\end{enumerate}

\subsection{Baselines}

\begin{itemize}[leftmargin=1.4em]
  \item \textbf{No-Interaction VLN:} ignores all post-instruction messages.
  \item \textbf{Dialogue Concatenation:} appends messages to the original
  instruction and replans without explicit task-state updates.
  \item \textbf{Always Ask:} requests help at every candidate ambiguity.
  \item \textbf{Threshold Ask:} asks when navigation entropy exceeds a tuned
  threshold.
  \item \textbf{Oracle Initiative:} receives oracle intervention timing or
  answers, providing an upper bound.
  \item \textbf{\method Baseline:} uses dynamic task memory and the learned
  initiative manager described in \cref{sec:method}.
\end{itemize}

\subsection{Implementation Details}

\tbd{Specify image encoder, language encoder, map representation, training
schedule, hardware, random seeds, and all cost coefficients.} We will report
mean and 95\% confidence intervals over at least three seeds. Counterfactual
branches will be evaluated with paired bootstrap tests. Model selection will
use validation-unseen data without access to test intervention paraphrases.

\subsection{Main Results}

\begin{table}[t]
\centering
\small
\resizebox{\linewidth}{!}{
\begin{tabular}{lrrrrrr}
\toprule
\textbf{Method} & \textbf{SR$\uparrow$} & \textbf{SPL$\uparrow$} &
\textbf{Recovery$\uparrow$} & \textbf{Viol.$\downarrow$} &
\textbf{Msgs.$\downarrow$} & \textbf{Utility$\uparrow$}\\
\midrule
No-Interaction VLN & TBD & TBD & TBD & TBD & 0 & TBD\\
Dialogue Concatenation & TBD & TBD & TBD & TBD & TBD & TBD\\
Always Ask & TBD & TBD & TBD & TBD & TBD & TBD\\
Threshold Ask & TBD & TBD & TBD & TBD & TBD & TBD\\
\method Baseline & TBD & TBD & TBD & TBD & TBD & TBD\\
Oracle Initiative & TBD & TBD & TBD & TBD & TBD & TBD\\
\bottomrule
\end{tabular}}
\caption{Main results on the test-unseen mixed-initiative track. No empirical
values are claimed in this draft.}
\label{tab:main}
\end{table}

\subsection{Initiative Direction and Timing}

We compare HI, RI, and MI tracks under matched episode prefixes. The analysis
will distinguish failures caused by navigation, task-state updating, question
selection, and unavailable users. Timing experiments evaluate early, middle,
and late interventions within each event window. We hypothesize that earlier
intervention reduces wasted motion, while a small amount of additional visual
evidence can make some later interventions easier to ground; this hypothesis
must be tested rather than assumed.

\begin{table}[t]
\centering
\small
\begin{tabular}{lrrrr}
\toprule
\textbf{Condition} & \textbf{SR} & \textbf{Recovery} &
\textbf{Extra Dist.} & \textbf{Human Cost}\\
\midrule
HI only & TBD & TBD & TBD & TBD\\
RI only & TBD & TBD & TBD & TBD\\
MI & TBD & TBD & TBD & TBD\\
MI + delayed user & TBD & TBD & TBD & TBD\\
MI + intermittent user & TBD & TBD & TBD & TBD\\
\bottomrule
\end{tabular}
\caption{Planned comparison of initiative and availability conditions.}
\label{tab:initiative}
\end{table}

\subsection{Ablations and Human Evaluation}

We ablate structured task memory, interaction history, availability modeling,
counterfactual training, and the communication-cost term. A human evaluation
will measure whether questions are timely and answerable, whether interventions
are correctly reflected in subsequent behavior, perceived workload, and trust.
Annotators will view paired rollouts with randomized method identity.

\section{Limitations and Broader Impact}
\label{sec:limitations}

Simulation cannot fully reproduce household noise, social expectations, sensor
failure, or the consequences of unsafe motion. The initial benchmark uses
short-horizon navigation actions and text interaction; speech recognition,
gesture, manipulation, and multi-person authority are outside its first scope.
Human messages may also encode cultural or household-specific assumptions.
Models trained on synthetic or templated language risk learning annotation
artifacts.

Mixed-initiative systems can reduce supervision burden, but poorly calibrated
asking may annoy users, and incorrect compliance with an intervention may
create safety risks. Benchmark success must therefore not be interpreted as
certification for deployment. Real-world testing should use conservative
motion limits, accessible stop controls, privacy-preserving data collection,
and explicit rules for resolving conflicting instructions.

\section{Conclusion}
\label{sec:conclusion}

We presented \method, a benchmark proposal for household VLN in which either
the robot or the human may initiate interaction. The formulation preserves
continuous navigation state, supports asynchronous human availability, and
represents language as updates to a dynamic task specification. The benchmark
combines controlled intervention windows, counterfactual branches, and metrics
that expose the trade-off between autonomy and human effort. The proposed
baseline and experiments provide a concrete path toward evaluating robots that
know when to ask and remain correctable whenever people choose to intervene.

\appendix

\section{Illustrative Episode Record}
\label{app:schema}

\begin{verbatim}
{
  "episode_id": "hssd_0123",
  "scene_id": "scene_x",
  "initial_instruction": "Go to the table in the living room.",
  "event": {
    "type": "reference_repair",
    "window": [24, 31],
    "trigger_step": 27,
    "message": "I meant the small table beside the sofa."
  },
  "availability": {
    "mode": "intermittent",
    "response_delay_steps": 4
  },
  "task_update": {
    "replace_goal_reference": "small_table_near_sofa"
  },
  "branches": [
    "no_intervention",
    "human_intervention",
    "oracle_intervention"
  ]
}
\end{verbatim}

\section{Additional Interaction Examples}

\paragraph{Human-initiated correction.}
The robot approaches the right-hand corridor. The user says, ``That route is
blocked; use the corridor behind you.'' The correct behavior is to update the
route constraint, preserve the destination, and replan from the current pose.

\paragraph{Robot-initiated clarification.}
At two visually similar doors, the robot asks, ``Do you mean the bedroom beside
the bathroom or the one at the end of the hall?'' If no answer arrives, the
agent must follow the configured conservative policy rather than assuming
continuous supervision.

\paragraph{Mixed initiative.}
The robot asks about an ambiguous destination. Before answering that question,
the user replaces the task: ``Never mind; return to the living room.'' The
system should invalidate the pending question and prioritize the new goal.

\section{Planned Release}

The intended release includes episode-generation code, simulator
configurations, state snapshots for counterfactual replay, train/validation/test
metadata, annotation guidelines, validated utterances, baseline models,
evaluation scripts, and model cards describing known limitations.

\begin{thebibliography}{99}

\bibitem{anderson2018r2r}
P. Anderson, Q. Wu, D. Teney, J. Bruce, M. Johnson, N. Sunderhauf,
I. Reid, S. Gould, and A. van den Hengel.
\newblock Vision-and-language navigation: Interpreting visually-grounded
navigation instructions in real environments.
\newblock In \emph{CVPR}, 2018.

\bibitem{beer2014levels}
J. M. Beer, A. D. Fisk, and W. A. Rogers.
\newblock Toward a framework for levels of robot autonomy in human-robot
interaction.
\newblock \emph{Journal of Human-Robot Interaction}, 2014.

\bibitem{broad2017language}
A. Broad, J. Murphey, and B. Argall.
\newblock Learning models for shared control of human-machine systems with
unknown dynamics.
\newblock \emph{Robotics: Science and Systems workshop}, 2017.

\bibitem{cui2023lilac}
Y. Cui, P. Karamcheti, A. Pal, N. Shivakumar, P. Liang, and D. Sadigh.
\newblock No, to the right: Online language corrections for robotic
manipulation via shared autonomy.
\newblock In \emph{HRI}, 2023.

\bibitem{khanna2024hssd}
M. Khanna et al.
\newblock HSSD: Habitat synthetic scenes dataset.
\newblock In \emph{CVPR}, 2024.

\bibitem{krantz2020beyond}
J. Krantz, E. Wijmans, A. Majumdar, D. Batra, and S. Lee.
\newblock Beyond the nav-graph: Vision-and-language navigation in continuous
environments.
\newblock In \emph{ECCV}, 2020.

\bibitem{krantz2023iterative}
J. Krantz, S. Banerjee, W. Zhu, J. Corso, P. Anderson, S. Lee, and J. Thomason.
\newblock Iterative vision-and-language navigation.
\newblock In \emph{CVPR}, 2023.

\bibitem{ku2020rxr}
A. Ku et al.
\newblock Room-across-room: Multilingual vision-and-language navigation with
dense spatiotemporal grounding.
\newblock In \emph{EMNLP}, 2020.

\bibitem{nguyen2019hanna}
K. Nguyen and H. Daume III.
\newblock Help, Anna! Visual navigation with natural multimodal assistance via
retrospective curiosity-encouraging imitation learning.
\newblock In \emph{EMNLP-IJCNLP}, 2019.

\bibitem{puig2024habitat3}
X. Puig et al.
\newblock Habitat 3.0: A co-habitat for humans, avatars and robots.
\newblock In \emph{ICLR}, 2024.

\bibitem{savva2019habitat}
M. Savva et al.
\newblock Habitat: A platform for embodied AI research.
\newblock In \emph{ICCV}, 2019.

\bibitem{sharma2022correcting}
P. Sharma, A. Torralba, and J. Andreas.
\newblock Skill induction and planning with latent language.
\newblock In \emph{ACL}, 2022.

\bibitem{taioli2025collaborative}
F. Taioli et al.
\newblock Collaborative object navigation using dialogue.
\newblock In \emph{ICCV}, 2025.

\bibitem{thomason2020cvdn}
J. Thomason, M. Murray, M. Cakmak, and L. Zettlemoyer.
\newblock Vision-and-dialog navigation.
\newblock In \emph{Conference on Robot Learning}, 2020.

\bibitem{zhu2020cmn}
Y. Zhu, F. Zhu, Z. Zhan, B. Lin, J. Jiao, X. Chang, and X. Liang.
\newblock Vision-dialog navigation by exploring cross-modal memory.
\newblock In \emph{CVPR}, 2020.

\bibitem{zhu2021scoa}
Y. Zhu et al.
\newblock Self-motivated communication agent for real-world vision-dialog
navigation.
\newblock In \emph{ICCV}, 2021.

\end{thebibliography}

\end{document}
